% ============================================================
\chapter{Joint Distributions and Covariance}
\label{ch:joint}
% ============================================================

Until now, we have studied random variables in isolation. But in the real world, quantities are rarely independent: height and weight are correlated, a stock's return depends on the overall market, temperature and atmospheric pressure evolve jointly. To understand these dependencies, one must move from the law of a single variable to the \emph{joint distribution} of several variables. Francis Galton, in the 1880s, was the first to study correlation systematically --- by measuring the heights of parents and their children. Karl Pearson then formalized the correlation coefficient that bears his name. This chapter develops the fundamental tools: joint distributions, marginals, conditionals, covariance, and correlation.

\begin{intuition}
This chapter studies the joint distributions of two (or more) random variables,
marginal and conditional distributions, covariance, and correlation. These tools
are essential for modelling dependence between random quantities.
\end{intuition}

% ------------------------------------------------------------
\section{Joint Distributions}
% ------------------------------------------------------------

\subsection{Discrete case}

\begin{definition}[Discrete joint distribution]
Let $(X, Y)$ be a pair of discrete random variables. The \textbf{joint distribution}
is the function
\[
    p_{X,Y}(x_i, y_j) = \PP(X = x_i, Y = y_j), \quad
    \text{with } \sum_{i,j} p_{X,Y}(x_i, y_j) = 1.
\]
\end{definition}

\subsection{Continuous case}

\begin{definition}[Joint density]
The pair $(X, Y)$ has a \textbf{joint density} $f_{X,Y}$ if, for every Borel set
$B \subseteq \R^2$:
\[
    \PP((X,Y) \in B) = \iint_B f_{X,Y}(x, y)\,dx\,dy,
\]
with $f_{X,Y}(x,y) \geq 0$ and $\iint_{\R^2} f_{X,Y}(x,y)\,dx\,dy = 1$.
\end{definition}

\begin{definition}[Joint CDF]
\[
    F_{X,Y}(x, y) = \PP(X \leq x,\; Y \leq y).
\]
In the continuous case: $f_{X,Y}(x,y) = \frac{\partial^2 F_{X,Y}}{\partial x\,\partial y}(x,y)$.
\end{definition}

% ------------------------------------------------------------
\section{Marginal Distributions}
% ------------------------------------------------------------

\begin{definition}[Marginals]
The \textbf{marginal distributions} are obtained by summing (discrete) or
integrating (continuous) over the other variable:
\begin{align*}
    &\text{Discrete:} & p_X(x_i) &= \sum_j p_{X,Y}(x_i, y_j), &
    p_Y(y_j) &= \sum_i p_{X,Y}(x_i, y_j). \\
    &\text{Continuous:} & f_X(x) &= \int_{-\infty}^{+\infty} f_{X,Y}(x,y)\,dy, &
    f_Y(y) &= \int_{-\infty}^{+\infty} f_{X,Y}(x,y)\,dx.
\end{align*}
\end{definition}

\begin{warning}[title=\textbf{The joint determines the marginals, but not vice versa}]
From $f_{X,Y}$ one can compute $f_X$ and $f_Y$. But knowing only $f_X$ and $f_Y$
is not enough to reconstruct $f_{X,Y}$ without information about the dependence
between $X$ and $Y$.
\end{warning}

\begin{example}
Let $(X,Y)$ have joint density $f_{X,Y}(x,y) = 6(1-y)$ for $0 < x < y < 1$
and $0$ otherwise. Then:
\[
    f_X(x) = \int_x^1 6(1-y)\,dy = 3(1-x)^2, \quad 0 < x < 1,
\]
\[
    f_Y(y) = \int_0^y 6(1-y)\,dx = 6y(1-y), \quad 0 < y < 1.
\]
\end{example}

% ------------------------------------------------------------
\section{Conditional Distributions}
% ------------------------------------------------------------

\begin{definition}[Conditional density]
The \textbf{conditional density} of $Y$ given $X = x$ is
\[
    f_{Y \mid X}(y \mid x) = \frac{f_{X,Y}(x, y)}{f_X(x)}, \quad f_X(x) > 0.
\]
Similarly in the discrete case:
$\PP(Y = y_j \mid X = x_i) = p_{X,Y}(x_i, y_j)/p_X(x_i)$.
\end{definition}

% ------------------------------------------------------------
\section{Independence}
% ------------------------------------------------------------

\begin{theorem}[Characterisation of independence]
$X$ and $Y$ are independent if and only if:
\begin{itemize}
    \item \textbf{Discrete:} $p_{X,Y}(x,y) = p_X(x)\,p_Y(y)$ for all $(x,y)$.
    \item \textbf{Continuous:} $f_{X,Y}(x,y) = f_X(x)\,f_Y(y)$ for all $(x,y)$.
    \item \textbf{Via the CDF:} $F_{X,Y}(x,y) = F_X(x)\,F_Y(y)$ for all $(x,y)$.
\end{itemize}
\end{theorem}

\begin{proposition}[Factorisation criterion]
The continuous pair $(X,Y)$ is independent if and only if there exist functions
$g, h$ such that $f_{X,Y}(x,y) = g(x)\,h(y)$ on the support (which must be a
Cartesian product).
\end{proposition}

% ------------------------------------------------------------
\section{Covariance and Correlation}
% ------------------------------------------------------------

\begin{definition}[Covariance]
\[
    \Cov(X, Y) = \E[(X - \E[X])(Y - \E[Y])] = \E[XY] - \E[X]\E[Y].
\]
\end{definition}

\begin{proposition}[Properties of covariance]
\begin{enumerate}[label=(\roman*)]
    \item $\Cov(X, X) = \Var(X)$.
    \item $\Cov(X, Y) = \Cov(Y, X)$ (symmetry).
    \item $\Cov(aX+b, cY+d) = ac\,\Cov(X,Y)$ (bilinearity).
    \item $\Var(X+Y) = \Var(X) + \Var(Y) + 2\Cov(X,Y)$.
    \item If $X \indep Y$, then $\Cov(X,Y) = 0$.
    \item $\Var\!\left(\sum_{i=1}^n X_i\right) = \sum_{i=1}^n \Var(X_i)
          + 2\sum_{i<j}\Cov(X_i, X_j)$.
\end{enumerate}
\end{proposition}

\begin{definition}[Correlation coefficient]
\[
    \rho(X, Y) = \frac{\Cov(X, Y)}{\sqrt{\Var(X)\,\Var(Y)}}.
\]
\end{definition}

\begin{theorem}[Properties of $\rho$]
\begin{enumerate}[label=(\roman*)]
    \item $-1 \leq \rho(X,Y) \leq 1$ \quad (consequence of Cauchy--Schwarz).
    \item $\abs{\rho(X,Y)} = 1$ iff $Y = aX + b$ a.s.\ with $a \neq 0$.
    \item $\rho(X,Y) = 0$ means $X$ and $Y$ are \textbf{uncorrelated}.
\end{enumerate}
\end{theorem}

\begin{warning}[title=\textbf{Uncorrelated $\neq$ independent}]
$\Cov(X,Y) = 0$ does not imply independence. Example: if $X \sim \mathcal{N}(0,1)$
and $Y = X^2$, then $\Cov(X,Y) = \E[X^3] = 0$ (by symmetry), yet $X$ and $Y$
are clearly not independent.
\end{warning}

% ------------------------------------------------------------
\section{Bivariate Normal Distribution}
% ------------------------------------------------------------

\begin{definition}[Bivariate normal]
The vector $(X, Y)$ has a \textbf{bivariate normal} distribution
$\mathcal{N}_2(\boldsymbol{\mu}, \Sigma)$ with
$\boldsymbol{\mu} = (\mu_X, \mu_Y)^T$ and
$\Sigma = \bigl(\begin{smallmatrix} \sigma_X^2 & \rho\sigma_X\sigma_Y \\
\rho\sigma_X\sigma_Y & \sigma_Y^2 \end{smallmatrix}\bigr)$ if
\[
    f_{X,Y}(x,y) = \frac{1}{2\pi\sigma_X\sigma_Y\sqrt{1-\rho^2}}
    \exp\!\left(-\frac{1}{2(1-\rho^2)}\left[
    \frac{(x-\mu_X)^2}{\sigma_X^2}
    - 2\rho\frac{(x-\mu_X)(y-\mu_Y)}{\sigma_X\sigma_Y}
    + \frac{(y-\mu_Y)^2}{\sigma_Y^2}
    \right]\right).
\]
\end{definition}

\begin{theorem}[Properties of the bivariate normal]
\begin{enumerate}[label=(\roman*)]
    \item The marginals are $X \sim \mathcal{N}(\mu_X, \sigma_X^2)$ and
          $Y \sim \mathcal{N}(\mu_Y, \sigma_Y^2)$.
    \item $\rho(X,Y) = \rho$ (the distribution parameter).
    \item $X$ and $Y$ are independent if and only if $\rho = 0$.
    \item Every linear combination $aX + bY$ is normal.
\end{enumerate}
\end{theorem}

\begin{remark}
Property (iii) is specific to the normal distribution: in general, uncorrelated
does not imply independent, but for a Gaussian vector it does.
\end{remark}

% ------------------------------------------------------------
\section{Sum of Random Variables}
% ------------------------------------------------------------

\begin{theorem}[Convolution]
If $X$ and $Y$ are independent with densities $f_X$ and $f_Y$, then
$Z = X + Y$ has density given by the \textbf{convolution}:
\[
    f_Z(z) = (f_X * f_Y)(z) = \int_{-\infty}^{+\infty} f_X(t)\,f_Y(z-t)\,dt.
\]
\end{theorem}

\begin{example}
If $X \sim \mathrm{Exp}(\lambda)$ and $Y \sim \mathrm{Exp}(\lambda)$ are independent,
then for $z > 0$:
\[
    f_Z(z) = \int_0^z \lambda e^{-\lambda t}\,\lambda e^{-\lambda(z-t)}\,dt
    = \lambda^2 e^{-\lambda z}\int_0^z dt = \lambda^2 z\,e^{-\lambda z}.
\]
This is the density of $\Gamma(2, \lambda)$.
\end{example}

% ------------------------------------------------------------
\section{Exercises}
% ------------------------------------------------------------

\begin{exercise}
Let $(X,Y)$ have density $f(x,y) = c\,xy$ for $0 < x < 1$, $0 < y < 1$
and $0$ otherwise. Find $c$, the marginals, and check independence.
\end{exercise}

\begin{exercise}
Let $(X,Y)$ have density $f(x,y) = 2$ for $0 < x < y < 1$, $0$ otherwise.
Compute $f_X$, $f_Y$, $\E[X]$, $\E[Y]$, $\Cov(X,Y)$, $\rho(X,Y)$.
\end{exercise}

\begin{exercise}
Show that if $(X,Y) \sim \mathcal{N}_2$ with $\rho = 0$, then $X$ and $Y$
are independent (factorise the joint density).
\end{exercise}

\begin{exercise}
Let $X \sim \mathcal{N}(0,1)$ and $Y = X^2$. Verify that $\Cov(X,Y) = 0$
but $X$ and $Y$ are not independent.
\end{exercise}

\begin{exercise}
Find by convolution the density of $Z = X + Y$ where $X, Y$ are i.i.d.\
$\mathcal{N}(0,1)$.
\end{exercise}

\begin{exercise}
Let $(X,Y)$ be uniform on the disc $\{(x,y) : x^2+y^2 \leq 1\}$.
Find the marginals, $\E[X]$, $\Var(X)$, $\Cov(X,Y)$.
\end{exercise}

\begin{exercise}
A point $(X,Y)$ is chosen uniformly in the triangle with vertices
$(0,0)$, $(1,0)$, $(0,1)$. Find the joint density, the marginals,
and compute $\rho(X,Y)$.
\end{exercise}

\begin{exercise}
Let $X_1, X_2, X_3$ be i.i.d.\ $\mathrm{Exp}(1)$. Compute
$\Cov(X_1+X_2,\; X_2+X_3)$.
\end{exercise}
